\documentclass[a4paper,12pt]{article}
\usepackage[utf8]{inputenc}
\usepackage[ngerman]{babel}
\usepackage{amsmath,amssymb}
\usepackage{geometry}\geometry{margin=2.5cm}
\usepackage{enumitem}
\usepackage{booktabs}
\usepackage{tikz}
\usepackage{pgfplots}\pgfplotsset{compat=1.18}
\usepackage{tcolorbox}
\tcbuselibrary{breakable}
\usepackage{xcolor}
\newtcolorbox{begriffe}[1][]{colback=yellow!10, colframe=yellow!60!black, fonttitle=\bfseries, title={Was bedeuten die Begriffe?}, breakable, #1}
\newtcolorbox{wastun}[1][]{colback=orange!5, colframe=orange!60!black, fonttitle=\bfseries, title={Was ist zu tun?}, breakable, #1}
\newcommand{\piv}[1]{{\color{red}\boxed{#1}}}

\title{\"Ubungsblatt 7 -- L\"osungen \\ \large Linear Algebra}
\author{}
\date{}

\begin{document}
\maketitle

%% ============================================================
\subsection*{Aufgabe 25 -- (nicht zur Abgabe)}
%% ============================================================

\"Ubungsaufgabe zum freien \"Uben mit zuf\"alligen $3\times 3$- und $4 \times 4$-Matrizen. Keine Abgabe.

\bigskip
\hrule
\bigskip

%% ============================================================
\subsection*{Aufgabe 26}
%% ============================================================

\textbf{Angabe.} Berechne die Determinanten effizient (per Hand, ohne Rechner):
\[
A = \begin{pmatrix} 234 & 235 & 236 \\ 235 & 236 & 237 \\ 236 & 237 & 239 \end{pmatrix},\qquad
B = \begin{pmatrix} 21 & 23 & 0 & 27 & 2 \\ 0 & 2 & 0 & 7 & 0 \\ 11 & 12 & 13 & 14 & 15 \\ 1 & 41 & 0 & 43 & 0 \\ 0 & 1 & 0 & 4 & 0 \end{pmatrix}.
\]

\begin{begriffe}
\begin{itemize}[leftmargin=*, nosep]
    \item \textbf{Determinante} -- eine Zahl, die einer quadratischen Matrix zugeordnet wird. $\det(A) = 0 \iff$ Zeilen (oder Spalten) sind linear abh\"angig.
    \item \textbf{Eigenschaft 1: Zeilenumformungen.} $Z_i \to Z_i + \lambda Z_j$ \"andert die Determinante \textbf{nicht}.
    \item \textbf{Eigenschaft 2: Zeilentausch.} \"Andert das Vorzeichen der Determinante.
    \item \textbf{Eigenschaft 3: Spaltenumformungen.} Genau wie bei Zeilen.
    \item \textbf{Laplace-Entwicklung} -- Verfahren, um eine gro\ss e Determinante in mehrere kleinere zu zerlegen. Anschaulich: ich w\"ahle eine Zeile (oder Spalte), gehe jeden Eintrag dieser Zeile durch und reduziere damit die $n \times n$-Determinante auf $n$ kleinere $(n-1) \times (n-1)$-Determinanten. Genaue Anleitung weiter unten beim ersten Einsatz.
    \item \textbf{Trick}: Zeilen/Spalten mit vielen Nullen suchen $\Rightarrow$ in der Summe werden viele Terme automatisch $0$, und die Rechnung wird viel k\"urzer.
\end{itemize}
\end{begriffe}

\begin{wastun}
\begin{enumerate}[leftmargin=*, nosep]
    \item \textbf{Matrix $A$:} Zahlen sehen wie aufeinanderfolgende Werte aus. Mit Zeilenumformungen $Z_2 \to Z_2 - Z_1$, $Z_3 \to Z_3 - Z_1$ entstehen kleine Zahlen.
    \item \textbf{Matrix $B$:} Die Spalte 3 hat fast nur Nullen (nur ein Eintrag $\neq 0$). Laplace-Entwicklung nach Spalte 3 reduziert auf eine $4 \times 4$-Matrix. Die wieder hat eine Spalte mit fast nur Nullen $\to$ weiter entwickeln, bis nur noch $2 \times 2$ \"ubrig.
\end{enumerate}
\end{wastun}

%% --- 26 A ---
\subsubsection*{Determinante von $A$}

\textbf{Schritt 1: Zeilenumformungen, um die gro\ss en Zahlen loszuwerden.}

\textbf{(a) $Z_2 \to Z_2 - Z_1$ und $Z_3 \to Z_3 - Z_1$} (Determinante \textbf{unver\"andert}):
\[
\det A = \det \begin{pmatrix} 234 & 235 & 236 \\ 1 & 1 & 1 \\ 2 & 2 & 3 \end{pmatrix}
\]

\textbf{(b) $Z_3 \to Z_3 - 2Z_2$} (Determinante \textbf{unver\"andert}):
\[
= \det \begin{pmatrix} 234 & 235 & 236 \\ 1 & 1 & 1 \\ 0 & 0 & 1 \end{pmatrix}
\]

\textbf{Schritt 2: Laplace-Entwicklung nach Zeile 3.}

Zeile 3 ist $(0, 0, 1)$. Die zwei Nullen liefern Beitrag $0$, also bleibt nur Eintrag $(3,3) = 1$. Streiche \textbf{Zeile 3 und Spalte 3} weg (grau):

\[
\begin{pmatrix}
234 & 235 & \color{gray}236 \\
1 & 1 & \color{gray}1 \\
\color{gray}0 & \color{gray}0 & \color{gray}1
\end{pmatrix}
\;\longrightarrow\;
\det = 1 \cdot \underbrace{(-1)^{3+3}}_{+1} \cdot \det \begin{pmatrix} 234 & 235 \\ 1 & 1 \end{pmatrix}
= 234 - 235 = -1.
\]

\[
\fbox{\parbox{0.93\textwidth}{\centering
$\det A = -1$.
}}
\]

%% --- 26 B ---
\subsubsection*{Determinante von $B$}

\textbf{Idee:} Spalten/Zeilen mit fast nur Nullen finden und wegklappen, bis nur noch $2 \times 2$ \"ubrig.

\textbf{Schritt 1: $5\times 5 \to 4\times 4$. Entwicklung nach Spalte 3.}

Spalte 3 ist $(0, 0, 13, 0, 0)$. Nur Eintrag $(3,3) = 13$ liefert Beitrag. Streiche Zeile 3 und Spalte 3:

\[
\begin{pmatrix}
21 & 23 & \color{gray}0 & 27 & 2 \\
0 & 2 & \color{gray}0 & 7 & 0 \\
\color{gray}11 & \color{gray}12 & \color{gray}13 & \color{gray}14 & \color{gray}15 \\
1 & 41 & \color{gray}0 & 43 & 0 \\
0 & 1 & \color{gray}0 & 4 & 0
\end{pmatrix}
\;\longrightarrow\;
\det B = 13 \cdot \underbrace{(-1)^{3+3}}_{+1} \cdot \det \underbrace{\begin{pmatrix} 21 & 23 & 27 & 2 \\ 0 & 2 & 7 & 0 \\ 1 & 41 & 43 & 0 \\ 0 & 1 & 4 & 0 \end{pmatrix}}_{=: M}
\]

\textbf{Schritt 2: $4\times 4 \to 3\times 3$. Entwicklung von $M$ nach Spalte 4.}

Spalte 4 von $M$ ist $(2, 0, 0, 0)$. Nur Eintrag $(1,4) = 2$ liefert Beitrag. Streiche Zeile 1 und Spalte 4:

\[
\begin{pmatrix}
\color{gray}21 & \color{gray}23 & \color{gray}27 & \color{gray}2 \\
0 & 2 & 7 & \color{gray}0 \\
1 & 41 & 43 & \color{gray}0 \\
0 & 1 & 4 & \color{gray}0
\end{pmatrix}
\;\longrightarrow\;
\det M = 2 \cdot \underbrace{(-1)^{1+4}}_{-1} \cdot \det \underbrace{\begin{pmatrix} 0 & 2 & 7 \\ 1 & 41 & 43 \\ 0 & 1 & 4 \end{pmatrix}}_{=: N}
= -2 \cdot \det N.
\]

\textbf{Schritt 3: $3\times 3 \to 2\times 2$. Entwicklung von $N$ nach Spalte 1.}

Spalte 1 von $N$ ist $(0, 1, 0)$. Nur Eintrag $(2,1) = 1$ liefert Beitrag. Streiche Zeile 2 und Spalte 1:

\[
\begin{pmatrix}
\color{gray}0 & 2 & 7 \\
\color{gray}1 & \color{gray}41 & \color{gray}43 \\
\color{gray}0 & 1 & 4
\end{pmatrix}
\;\longrightarrow\;
\det N = 1 \cdot \underbrace{(-1)^{2+1}}_{-1} \cdot \det \begin{pmatrix} 2 & 7 \\ 1 & 4 \end{pmatrix}
= -1 \cdot (8 - 7) = -1.
\]

\textbf{Schritt 4: Zur\"uckrechnen.}

\[
\det M = -2 \cdot (-1) = 2, \qquad \det B = 13 \cdot 2 = 26.
\]

\[
\fbox{\parbox{0.93\textwidth}{\centering
$\det B = 26$.
}}
\]

\bigskip
\hrule
\bigskip

%% ============================================================
\subsection*{Aufgabe 27}
%% ============================================================

\textbf{Angabe.}
\[
A = \begin{pmatrix} a & 1 & 2 & 3 \\ 1 & 2 & 3 & 2 \\ 2 & 3 & 2 & 1 \\ 3 & 2 & 1 & b \end{pmatrix} \in \mathbb{R}^{4\times 4},\quad a, b \in \mathbb{R}.
\]
(a) $\det A$ berechnen. \quad (b) Bedingung an $a, b$, sodass $Ax = 0$ unendlich viele L\"osungen hat.

\begin{begriffe}
\begin{itemize}[leftmargin=*, nosep]
    \item \textbf{Parameter $a, b$} -- unbestimmte Zahlen, die im Ergebnis stehen bleiben d\"urfen. $\det A$ wird ein Ausdruck in $a$ und $b$.
    \item \textbf{$Ax = 0$ hat unendlich viele L\"osungen} $\iff$ $A$ ist nicht invertierbar $\iff$ $\det A = 0$. (Das homogene System hat \emph{immer} die triviale L\"osung $x = 0$; weitere L\"osungen existieren genau dann, wenn $\det A = 0$.)
\end{itemize}
\end{begriffe}

\begin{wastun}
\begin{enumerate}[leftmargin=*, nosep]
    \item \textbf{(a)} Determinante per Laplace-Entwicklung nach der ersten Zeile berechnen (4 Minoren, jeder $3\times 3$). Parameter $a, b$ bleiben mitgeschleppt.
    \item \textbf{(b)} $\det A = 0$ setzen und nach $a, b$ aufl\"osen.
\end{enumerate}
\end{wastun}

%% --- 27(a) ---
\subsubsection*{(a) Determinante von $A$}

\textbf{Schritt 1: Laplace-Entwicklung nach der ersten Zeile.}

Notation: $a_{ij}$ steht f\"ur den Eintrag in \textbf{Zeile $i$, Spalte $j$}. Hier:
\[
A = \begin{pmatrix}
\underbrace{a}_{a_{11}} & \underbrace{1}_{a_{12}} & \underbrace{2}_{a_{13}} & \underbrace{3}_{a_{14}} \\
1 & 2 & 3 & 2 \\
2 & 3 & 2 & 1 \\
3 & 2 & 1 & b
\end{pmatrix}
\]
also $a_{11} = a$, $a_{12} = 1$, $a_{13} = 2$, $a_{14} = 3$.

Vorzeichen f\"ur Zeile 1 (Schachbrett): $+\;-\;+\;-$. Damit:
\[
\det A = a \cdot M_{11} - 1 \cdot M_{12} + 2 \cdot M_{13} - 3 \cdot M_{14}.
\]
$M_{1j}$: Zeile 1 und Spalte $j$ wegstreichen.

\textbf{Schritt 2: Vier Minoren ausrechnen.}

Pro Minor $M_{1j}$ wird in $A$ \textbf{Zeile 1 und Spalte $j$} weggestrichen (grau dargestellt) -- die schwarzen Eintr\"age bilden die $3\times 3$-Matrix.

\textbf{(a) $M_{11}$ -- Zeile 1 und Spalte 1 weg.}
\[
\begin{pmatrix}
\color{gray}a & \color{gray}1 & \color{gray}2 & \color{gray}3 \\
\color{gray}1 & 2 & 3 & 2 \\
\color{gray}2 & 3 & 2 & 1 \\
\color{gray}3 & 2 & 1 & b
\end{pmatrix}
\;\longrightarrow\;
M_{11} = \det \begin{pmatrix} 2 & 3 & 2 \\ 3 & 2 & 1 \\ 2 & 1 & b \end{pmatrix}
= 2(2b - 1) - 3(3b - 2) + 2(3 - 4)
= -5b + 2.
\]

\textbf{(b) $M_{12}$ -- Zeile 1 und Spalte 2 weg.}
\[
\begin{pmatrix}
\color{gray}a & \color{gray}1 & \color{gray}2 & \color{gray}3 \\
1 & \color{gray}2 & 3 & 2 \\
2 & \color{gray}3 & 2 & 1 \\
3 & \color{gray}2 & 1 & b
\end{pmatrix}
\;\longrightarrow\;
M_{12} = \det \begin{pmatrix} 1 & 3 & 2 \\ 2 & 2 & 1 \\ 3 & 1 & b \end{pmatrix}
= 1(2b - 1) - 3(2b - 3) + 2(2 - 6)
= -4b.
\]

\textbf{(c) $M_{13}$ -- Zeile 1 und Spalte 3 weg.}
\[
\begin{pmatrix}
\color{gray}a & \color{gray}1 & \color{gray}2 & \color{gray}3 \\
1 & 2 & \color{gray}3 & 2 \\
2 & 3 & \color{gray}2 & 1 \\
3 & 2 & \color{gray}1 & b
\end{pmatrix}
\;\longrightarrow\;
M_{13} = \det \begin{pmatrix} 1 & 2 & 2 \\ 2 & 3 & 1 \\ 3 & 2 & b \end{pmatrix}
= 1(3b - 2) - 2(2b - 3) + 2(4 - 9)
= -b - 6.
\]

\textbf{(d) $M_{14}$ -- Zeile 1 und Spalte 4 weg.}
\[
\begin{pmatrix}
\color{gray}a & \color{gray}1 & \color{gray}2 & \color{gray}3 \\
1 & 2 & 3 & \color{gray}2 \\
2 & 3 & 2 & \color{gray}1 \\
3 & 2 & 1 & \color{gray}b
\end{pmatrix}
\;\longrightarrow\;
M_{14} = \det \begin{pmatrix} 1 & 2 & 3 \\ 2 & 3 & 2 \\ 3 & 2 & 1 \end{pmatrix}
= 1(3 - 4) - 2(2 - 6) + 3(4 - 9)
= -8.
\]

\textbf{Schritt 3: Zusammensetzen.}

\begin{align*}
\det A &= a \cdot (-5b+2) - 1 \cdot (-4b) + 2 \cdot (-b - 6) - 3 \cdot (-8) \\
&= -5ab + 2a + 4b - 2b - 12 + 24 \\
&= -5ab + 2a + 2b + 12.
\end{align*}

\[
\fbox{\parbox{0.93\textwidth}{\centering
$\det A = -5ab + 2a + 2b + 12$.
}}
\]

%% --- 27(b) ---
\subsubsection*{(b) Bedingung f\"ur unendlich viele L\"osungen}

\textbf{Schritt 1: $\det A = 0$ ansetzen.}

$Ax = 0$ hat unendlich viele L\"osungen $\iff \det A = 0$:
\[
-5ab + 2a + 2b + 12 = 0 \quad\iff\quad 5ab - 2a - 2b - 12 = 0.
\]

\textbf{Schritt 2: Faktorisieren.}

\textbf{Warum \"uberhaupt faktorisieren?} Die Gleichung $5ab - 2a - 2b - 12 = 0$ steht in einer ``unsch\"onen'' Form -- man sieht $a$, $b$ und das Produkt $ab$ nebeneinander. Eine Form $(\dots)(\dots) = \text{Zahl}$ ist deutlich angenehmer: man hat dann zwei klare Faktoren, kann konkrete L\"osungspaare leichter bauen und sieht direkt, dass es \emph{unendlich} viele $(a, b)$ gibt.

\textbf{Warum mit $5$ multiplizieren?} Wir wollen die Gleichung in die Form $(5a + p)(5b + q) = \text{Zahl}$ bringen, weil das den Faktor $5ab$ aus der Originalgleichung mit dem Faktor $25ab = (5a)(5b)$ aus dem Produkt verbindet. Genauer:
\[
(5a + p)(5b + q) = 25ab + 5qa + 5pb + pq.
\]
Damit Koeffizientenvergleich m\"oglich wird, brauchen wir auf der linken Seite den Vorfaktor $25$ vor $ab$. Daher multiplizieren wir die ganze Gleichung mit $5$ -- das ist der Faktor, der $5ab$ zu $25ab$ macht, ohne neue Variablen einzuf\"uhren.

\textbf{Schritt 2a: Mit $5$ multiplizieren.}
\[
5 \cdot (5ab - 2a - 2b - 12) = 0 \;\Longleftrightarrow\; 25ab - 10a - 10b - 60 = 0.
\]

\textbf{Schritt 2b: Form $(5a + p)(5b + q)$ raten.} Die Koeffizienten von $a$ und $b$ in $25ab - 10a - 10b - 60$ sind beide $-10$. Aus
\[
(5a + p)(5b + q) = 25ab + 5qa + 5pb + pq
\]
folgt $5q = -10$ und $5p = -10$, also $p = q = -2$. Damit:
\[
(5a - 2)(5b - 2) = 25ab - 10a - 10b + 4.
\]

\textbf{Schritt 2c: Konstanten ausgleichen.} In unserer Gleichung steht $-60$ statt $+4$. Differenz: $4 - (-60) = 64$. Genau diese $64$ landet auf der rechten Seite:
\[
25ab - 10a - 10b - 60 = 0
\;\Longleftrightarrow\; \underbrace{(25ab - 10a - 10b + 4)}_{= (5a-2)(5b-2)} - 64 = 0
\;\Longleftrightarrow\; (5a-2)(5b-2) = 64.
\]

\[
\fbox{\parbox{0.93\textwidth}{\centering
$Ax = 0$ hat unendlich viele L\"osungen $\iff$ $(5a-2)(5b-2) = 64$.\\[2pt]
\"Aquivalent: $5ab - 2a - 2b = 12$.
}}
\]

\bigskip
\hrule
\bigskip

%% ============================================================
\subsection*{Aufgabe 28}
%% ============================================================

\textbf{Angabe.} Sei $n \geq 3$ ungerade, $A \in \mathbb{R}^{n \times n}$ schiefsymmetrisch ($A^\top = -A$). Zeige: die Zeilen von $A$ sind linear abh\"angig.

\begin{begriffe}
\begin{itemize}[leftmargin=*, nosep]
    \item \textbf{Schiefsymmetrisch}: $A^\top = -A$, d.\,h.\ $a_{ij} = -a_{ji}$.
    \item Zeilen sind \textbf{linear abh\"angig} $\iff$ $\det A = 0$.
    \item Es gen\"ugt also, $\det A = 0$ zu zeigen.
\end{itemize}
\end{begriffe}

\begin{wastun}
\begin{enumerate}[leftmargin=*, nosep]
    \item Mit Determinanten-Eigenschaften zeigen, dass $\det A = -\det A$ folgt.
    \item Daraus $\det A = 0$, daher Zeilen linear abh\"angig.
\end{enumerate}
\end{wastun}

\textbf{Schritt 1: Zwei Determinanten-Identit\"aten.}

\textbf{(a) Transponieren \"andert die Determinante nicht:}
\[
\det(A^\top) = \det(A).
\]

\textbf{(b) Multiplikation mit Skalar zieht $\lambda^n$ heraus:} F\"ur $\lambda \in \mathbb{R}$ gilt $\det(\lambda A) = \lambda^n \det(A)$. Insbesondere mit $\lambda = -1$:
\[
\det(-A) = (-1)^n \det(A).
\]

\textbf{Schritt 2: Beides verkn\"upfen.}

Wegen $A^\top = -A$:
\[
\det(A) = \det(A^\top) = \det(-A) = (-1)^n \det(A).
\]

\textbf{Schritt 3: $n$ ungerade ausnutzen.}

F\"ur ungerades $n$ ist $(-1)^n = -1$, also:
\[
\det(A) = -\det(A) \;\Rightarrow\; 2 \det(A) = 0 \;\Rightarrow\; \det(A) = 0.
\]

\textbf{Schritt 4: Folgerung.}

$\det(A) = 0 \Rightarrow$ die Zeilen von $A$ sind linear abh\"angig. $\square$

\bigskip
\hrule
\bigskip

%% ============================================================
\subsection*{Aufgabe 29}
%% ============================================================

\textbf{Angabe.} Sei $n \geq 3$, $A \in \mathbb{R}^{n \times n}$, jede Zeile von $A$ bildet eine arithmetische Folge. Zeige: $\det(A) = 0$.

\begin{begriffe}
\begin{itemize}[leftmargin=*, nosep]
    \item \textbf{Arithmetische Folge}: aufeinanderfolgende Differenzen sind konstant. Zeile $i$ hat Form $(a_i,\; a_i + d_i,\; a_i + 2d_i,\; \dots,\; a_i + (n-1)d_i)$ mit Anfangswert $a_i$ und Differenz $d_i$.
    \item \textbf{Charakterisierung}: Drei aufeinanderfolgende Glieder $x, y, z$ einer arithmetischen Folge erf\"ullen $z - 2y + x = 0$.
\end{itemize}
\end{begriffe}

\subsubsection*{Argument 1: Linear abh\"angige Spalten}

\textbf{Schritt 1: Spaltenrelation aufstellen.}

Sei $S_j$ die $j$-te Spalte von $A$, also $S_j = (a_1 + (j-1)d_1, \, a_2 + (j-1)d_2, \, \dots, \, a_n + (j-1)d_n)^\top$.

Mit $u = (a_1, \dots, a_n)^\top$ und $v = (d_1, \dots, d_n)^\top$ gilt:
\[
S_j = u + (j-1) v.
\]

Insbesondere: $S_1 = u$, $\quad S_2 = u + v$, $\quad S_3 = u + 2v$.

\textbf{Schritt 2: Lineare Abh\"angigkeit von $S_1, S_2, S_3$.}

\[
S_3 - 2 S_2 + S_1 = (u+2v) - 2(u+v) + u = 0.
\]

Das ist eine \textbf{nichttriviale} Linearkombination der Null mit Koeffizienten $(1, -2, 1)$.

\textbf{Schritt 3: Folgerung.}

Drei Spalten sind linear abh\"angig $\Rightarrow$ alle Spalten zusammen sind linear abh\"angig $\Rightarrow$ $\det(A) = 0$. $\square$

\subsubsection*{Argument 2: Spaltenumformung gibt Nullspalte}

\textbf{Schritt 1: Spaltenumformung anwenden.}

Determinante \"andert sich nicht durch $S_3 \to S_3 - 2 S_2 + S_1$. F\"ur den Eintrag in Zeile $i$ der neuen Spalte 3:
\[
(a_i + 2d_i) - 2(a_i + d_i) + (a_i) = a_i(1 - 2 + 1) + d_i(2 - 2 + 0) = 0.
\]

\textbf{Schritt 2: Folgerung.}

Die neue Spalte 3 ist die Nullspalte. Eine Matrix mit einer Nullspalte hat Determinante $0$. $\square$

\subsubsection*{Argument 3: $Ax = 0$ hat nichttriviale L\"osung}

W\"ahle $x = (1, -2, 1, 0, \dots, 0)^\top \in \mathbb{R}^n$. F\"ur jede Zeile $i$:
\[
(Ax)_i = 1 \cdot a_i + (-2)(a_i + d_i) + 1 \cdot (a_i + 2d_i) = 0.
\]

Also $Ax = 0$ mit $x \neq 0$ $\Rightarrow$ Spalten linear abh\"angig $\Rightarrow$ $\det(A) = 0$. $\square$

\subsubsection*{Argument 4: Dimension des Spaltenraums}

Aus Schritt 1 von Argument 1: jede Spalte $S_j = u + (j-1) v$ liegt in $\text{span}(u, v)$. Also hat der Spaltenraum H\"ochstdimension $2$. F\"ur $n \geq 3$ ist $\text{rang}(A) \leq 2 < n$, daher $\det(A) = 0$. $\square$

\bigskip
\hrule
\bigskip

%% ============================================================
\subsection*{Aufgabe 30}
%% ============================================================

\textbf{Angabe.}
\[
\begin{cases} 2x_1 + 5x_2 + 3x_3 = 5 \\ -x_1 - 4x_2 + x_3 = -3 \\ 2x_2 - 5x_3 = -1 \end{cases}
\]
auf zwei Wegen l\"osen: (a) \"uber die adjungierte Matrix, (b) Cramer'sche Regel.

\[
A = \begin{pmatrix} 2 & 5 & 3 \\ -1 & -4 & 1 \\ 0 & 2 & -5 \end{pmatrix},\quad
b = \begin{pmatrix} 5 \\ -3 \\ -1 \end{pmatrix}.
\]

\begin{begriffe}
\begin{itemize}[leftmargin=*, nosep]
    \item \textbf{Kofaktor} $C_{ij} = (-1)^{i+j} M_{ij}$, wobei $M_{ij}$ die Minor-Determinante (Matrix ohne Zeile $i$ und Spalte $j$) ist.
    \item \textbf{Kofaktormatrix} $\text{cof}(A) = (C_{ij})$ -- die Matrix aller Kofaktoren.
    \item \textbf{Adjungierte} (Adjunkte) $\text{adj}(A) = \text{cof}(A)^\top$ -- transponierte Kofaktormatrix.
    \item \textbf{Inverse \"uber Adjunkte}: $A^{-1} = \tfrac{1}{\det A} \cdot \text{adj}(A)$ (falls $\det A \neq 0$).
    \item \textbf{Cramer'sche Regel}: $x_i = \dfrac{\det A_i}{\det A}$, wobei $A_i$ aus $A$ entsteht, indem man die $i$-te Spalte durch $b$ ersetzt.
\end{itemize}
\end{begriffe}

\subsubsection*{Schritt 1: $\det A$ berechnen (f\"ur beide Methoden gebraucht)}

Laplace-Entwicklung nach erster Spalte (oder direkt Sarrus):
\begin{align*}
\det A &= 2 \cdot \det \begin{pmatrix} -4 & 1 \\ 2 & -5 \end{pmatrix} - 5 \cdot \det \begin{pmatrix} -1 & 1 \\ 0 & -5 \end{pmatrix} + 3 \cdot \det \begin{pmatrix} -1 & -4 \\ 0 & 2 \end{pmatrix} \\
&= 2 (20 - 2) - 5 (5 - 0) + 3 (-2 - 0) \\
&= 36 - 25 - 6 = 5.
\end{align*}

$\det A = 5 \neq 0 \Rightarrow$ System eindeutig l\"osbar.

%% --- 30(a) Adjugate ---
\subsubsection*{(a) L\"osung \"uber die adjungierte Matrix}

\textbf{Schritt 2: Alle 9 Kofaktoren $C_{ij} = (-1)^{i+j} M_{ij}$ berechnen.}

\begin{itemize}[leftmargin=*, nosep]
\item $C_{11} = +\det \begin{pmatrix} -4 & 1 \\ 2 & -5 \end{pmatrix} = 20 - 2 = 18$
\item $C_{12} = -\det \begin{pmatrix} -1 & 1 \\ 0 & -5 \end{pmatrix} = -(5 - 0) = -5$
\item $C_{13} = +\det \begin{pmatrix} -1 & -4 \\ 0 & 2 \end{pmatrix} = -2 - 0 = -2$
\item $C_{21} = -\det \begin{pmatrix} 5 & 3 \\ 2 & -5 \end{pmatrix} = -(-25 - 6) = 31$
\item $C_{22} = +\det \begin{pmatrix} 2 & 3 \\ 0 & -5 \end{pmatrix} = -10 - 0 = -10$
\item $C_{23} = -\det \begin{pmatrix} 2 & 5 \\ 0 & 2 \end{pmatrix} = -(4 - 0) = -4$
\item $C_{31} = +\det \begin{pmatrix} 5 & 3 \\ -4 & 1 \end{pmatrix} = 5 - (-12) = 17$
\item $C_{32} = -\det \begin{pmatrix} 2 & 3 \\ -1 & 1 \end{pmatrix} = -(2 - (-3)) = -5$
\item $C_{33} = +\det \begin{pmatrix} 2 & 5 \\ -1 & -4 \end{pmatrix} = -8 - (-5) = -3$
\end{itemize}

\textbf{Schritt 3: Kofaktormatrix und Adjunkte (transponieren).}

\[
\text{cof}(A) = \begin{pmatrix} 18 & -5 & -2 \\ 31 & -10 & -4 \\ 17 & -5 & -3 \end{pmatrix},\qquad
\text{adj}(A) = \text{cof}(A)^\top = \begin{pmatrix} 18 & 31 & 17 \\ -5 & -10 & -5 \\ -2 & -4 & -3 \end{pmatrix}.
\]

\textbf{Schritt 4: $A^{-1} = \tfrac{1}{\det A} \cdot \text{adj}(A)$ und $x = A^{-1} b$.}

\begin{align*}
A^{-1} b &= \frac{1}{5} \begin{pmatrix} 18 & 31 & 17 \\ -5 & -10 & -5 \\ -2 & -4 & -3 \end{pmatrix} \begin{pmatrix} 5 \\ -3 \\ -1 \end{pmatrix} \\
&= \frac{1}{5} \begin{pmatrix} 18 \cdot 5 + 31 \cdot (-3) + 17 \cdot (-1) \\ -5 \cdot 5 + (-10)(-3) + (-5)(-1) \\ -2 \cdot 5 + (-4)(-3) + (-3)(-1) \end{pmatrix}
= \frac{1}{5} \begin{pmatrix} 90 - 93 - 17 \\ -25 + 30 + 5 \\ -10 + 12 + 3 \end{pmatrix} \\
&= \frac{1}{5} \begin{pmatrix} -20 \\ 10 \\ 5 \end{pmatrix}
= \begin{pmatrix} -4 \\ 2 \\ 1 \end{pmatrix}.
\end{align*}

\[
\fbox{\parbox{0.93\textwidth}{\centering
$x = (x_1, x_2, x_3) = (-4,\; 2,\; 1)$.
}}
\]

%% --- 30(b) Cramer ---
\subsubsection*{(b) L\"osung mit der Cramer'schen Regel}

\textbf{Schritt 5: Drei Hilfsmatrizen $A_1, A_2, A_3$ aufstellen.}

$A_i$ entsteht aus $A$, indem die $i$-te Spalte durch $b$ ersetzt wird.

\textbf{(a) $A_1$ -- Spalte 1 wird $b$:}
\[
A_1 = \begin{pmatrix} 5 & 5 & 3 \\ -3 & -4 & 1 \\ -1 & 2 & -5 \end{pmatrix}, \quad
\det A_1 = 5(20 - 2) - 5(15 + 1) + 3(-6 - 4) = 90 - 80 - 30 = -20.
\]

\textbf{(b) $A_2$ -- Spalte 2 wird $b$:}
\[
A_2 = \begin{pmatrix} 2 & 5 & 3 \\ -1 & -3 & 1 \\ 0 & -1 & -5 \end{pmatrix}, \quad
\det A_2 = 2(15 + 1) - 5(5 - 0) + 3(1 - 0) = 32 - 25 + 3 = 10.
\]

\textbf{(c) $A_3$ -- Spalte 3 wird $b$:}
\[
A_3 = \begin{pmatrix} 2 & 5 & 5 \\ -1 & -4 & -3 \\ 0 & 2 & -1 \end{pmatrix}, \quad
\det A_3 = 2(4 + 6) - 5(1 - 0) + 5(-2 - 0) = 20 - 5 - 10 = 5.
\]

\textbf{Schritt 6: $x_i = \det A_i / \det A$.}

\[
x_1 = \frac{\det A_1}{\det A} = \frac{-20}{5} = -4, \quad
x_2 = \frac{\det A_2}{\det A} = \frac{10}{5} = 2, \quad
x_3 = \frac{\det A_3}{\det A} = \frac{5}{5} = 1.
\]

\[
\fbox{\parbox{0.93\textwidth}{\centering
$x = (-4,\; 2,\; 1)$ -- gleich wie in (a). $\checkmark$
}}
\]

\textbf{Probe:}
\begin{align*}
2(-4) + 5(2) + 3(1) &= -8 + 10 + 3 = 5 \;\checkmark \\
-(-4) - 4(2) + 1 &= 4 - 8 + 1 = -3 \;\checkmark \\
2(2) - 5(1) &= 4 - 5 = -1 \;\checkmark
\end{align*}

\end{document}
